\documentclass[11pt]{article}

\input{../preamble}

\usepackage{enumitem}
\newcounter{probnumcount}
\setcounter{probnumcount}{1}
\newcommand{\probnum}{\arabic{probnumcount}. \addtocounter{probnumcount}{1}}


\begin{document}
	
\begin{center}
\Large{\textbf{CS 330, Fall 2026, Homework 5 \\
Due Wednesday, October 7, 2026 \\
}}
\end{center}

\emph{For full credit, you must follow the guidelines laid out on the first page of homework 1.}\\


\subsection*{\probnum Job Fair Speedrun}

Boston University is hosting a massive Job Fair, and you are on a mission: hand out your CV and do an awkward-but-professional handshake with \emph{every} employer. Unfortunately, every employer at the job fair has taken CS 330, so they are terrifyingly efficient with their time management.

Employer $i\in\{1,\dots,n\}$ will only be at their booth during the time interval $[s_i,f_i]$ (continuous time). If you show up at any time $t\in[s_i,f_i]$, you successfully shake hands with employer $i$ (and possibly receive a sticker or a tote bag).

You want to find an \textit{acceptable} set of visit times $\{t_1,\dots,t_k\}$ such that each employer's interval $[s_i,f_i]$ contains at least one $t_j$ (i.e., $s_i \leq t_j \leq f_i$). Your goal is to minimize $k$, the number of visits to the job fair.

In the example of Figure~\ref{fig:visits}, the black vertical lines mark a set of acceptable visit times. Removing any one of them would make the set not acceptable. (But it is still not the smallest acceptable set.) 

\begin{figure}[h]
    \begin{center}
        \includegraphics[height=1in]{hw4-1.pdf}
    \end{center}
    \caption{An acceptable set of four visit times.}\label{fig:visits}
\end{figure}

\begin{enumerate}[label=(\alph*)]
\item (Do not hand in) Find a smallest acceptable set of attendance times for the example in Figure~\ref{fig:visits}.



\item Your friend suggests a greedy strategy: always pick the time that covers the largest number of remaining employers. Give an example where this fails to find the minimum hitting set.


\item Another friend observes: “If the input contains $k$ mutually disjoint intervals, then every acceptable set of visit times must have size at least $k$.” Is she right? Prove or disprove.



\item Design a polynomial-time algorithm that, given $[s_1,f_1],\dots,[s_n,f_n]$, outputs a minimum-size set of visit times. [\textit{Hint:} modify the interval scheduling algorithm.]



\item What is the running time of your algorithm in term of $n$?


\item Prove your algorithm is correct. [\textit{Hint:} Use the lower-bound observation from part (c).]



\end{enumerate}


\subsection*{\probnum Earliest Interval}
Suppose you are given a set of closed intervals $[s_i,f_i]$ for $i=1,\dots,n$. For simplicity, assume that all start and finish times are distinct (that is, there are no repetitions in the list $T=\{s_1,\dots,s_n,f_1,\dots,f_n\}$). For each finish time, we want to find the earliest-starting interval that contains that finish time. More concretely, given $f_k$, find an index $j$ such that
\[
f_k\in[s_j,f_j],
\]
and among all such intervals, $s_j$ is as small as possible.
Design an algorithm that outputs such an index for each finish time. The output may be given as $A$, a list of $n$ indices such that for the $k$th finish time, $A[k]$ is the value of the index you are asked to find. 




\end{document}